% ================================================================
% CHAPTER 5: ORTHOGONALITY AND THE STRUCTURE OF UNRELATED IDEAS
% ================================================================
\chapter{Orthogonality and the Structure of Unrelated Ideas}
\label{ch:orthogonality}

\epigraph{The world is everything that is the case.}{Ludwig Wittgenstein,
\textit{Tractatus Logico-Philosophicus}, 1}

\section{What Does It Mean for Ideas to Be Unrelated?}

In ordinary discourse, we say that two ideas are ``unrelated'' when they have
nothing to do with each other: chess strategy and molecular biology, Medieval
French poetry and semiconductor physics, the rules of cricket and the
philosophy of Spinoza. But ``unrelated'' is an imprecise term. Do unrelated
ideas merely fail to resonate, or do they have a stronger property?

In the Hilbert space model, ``unrelated'' has a precise meaning:
\emph{orthogonality}. Two ideas are orthogonal when their resonance is exactly
zero. This is stronger than merely ``low'' resonance: it means that the ideas
occupy genuinely independent dimensions of meaning space.

\begin{definition}[Orthogonal Ideas]
\label{def:orthogonal}
Two ideas $a, b \in \IS$ are \textbf{orthogonal}, written $a \perp b$, if:
\[
  \rs{a}{b} = \ip{\embed{a}}{\embed{b}} = 0.
\]
\end{definition}

\begin{definition}[Opposing Ideas]
\label{def:opposing}
Two ideas $a, b \in \IS$ are \textbf{opposing} if $\rs{a}{b} < 0$, and
\textbf{perfectly opposing} if $\hat{r}(a, b) = -1$ (i.e., $\embed{a}$ and
$\embed{b}$ are negative scalar multiples of each other).
\end{definition}

\begin{remark}[Orthogonality $\neq$ Opposition]
\label{rmk:orth-vs-opp}
This distinction is crucial and commonly confused. In ordinary language,
``unrelated'' and ``opposed'' are sometimes treated as the same thing ---
both are ``not agreeing.'' But in the geometry of meaning, they are radically
different:
\begin{itemize}
\item Orthogonal ideas ($\rs{a}{b} = 0$) are independent: knowing one tells
you nothing about the other. They compose without interference.
\item Opposing ideas ($\rs{a}{b} < 0$) are actively in conflict: each
\emph{diminishes} the other. They compose with cancellation.
\end{itemize}
``Capitalism'' and ``socialism'' are \emph{opposing} (negative resonance), not
orthogonal. ``Capitalism'' and ``photosynthesis'' are \emph{orthogonal} (zero
resonance), not opposing.
\end{remark}


\section{The Pythagorean Theorem of Ideas}

The central theorem of this chapter is the Pythagorean theorem for orthogonal
ideas: a result so fundamental that it deserves to be called the
\emph{Pythagorean theorem of ideas}.

\begin{theorem}[Pythagorean Theorem of Ideas]
\label{thm:pythagoras}
If $a \perp b$ (i.e., $\rs{a}{b} = 0$), then:
\[
  w(a \compose b)^2 = w(a)^2 + w(b)^2.
\]
In other words, $\nrm{\embed{a \compose b}}^2 = \nrm{\embed{a}}^2 +
\nrm{\embed{b}}^2$.
\end{theorem}

\begin{proof}
From Theorem~\ref{thm:norm-expansion} (Chapter~\ref{ch:depth-norm}):
\[
  w(a \compose b)^2 = w(a)^2 + w(b)^2 + 2\rs{a}{b} = w(a)^2 + w(b)^2 + 0
  = w(a)^2 + w(b)^2.
\]

In Lean~4:
\begin{lstlisting}
theorem pythagoras_ideas (h : rs a b = 0) :
    ‖φ (a ∘ b)‖ ^ 2 = ‖φ a‖ ^ 2 + ‖φ b‖ ^ 2 := by
  simp [embed_compose, norm_add_sq_real, rs] at *
  linarith
\end{lstlisting}
\end{proof}

\begin{interpretation}[Orthogonal Ideas Compose Without Interference]
\label{interp:pythagoras}
The Pythagorean theorem of ideas is a precise statement of a profound
intuition: \emph{unrelated ideas compose without interference}. When two
ideas occupy completely independent dimensions of meaning space, combining
them produces a composite whose weight is determined by the Pythagorean
formula --- exactly as one would expect from two independent contributions
that neither reinforce nor cancel each other.

This has important implications for intellectual life. When a scholar combines
expertise in two orthogonal domains --- say, medieval history and computational
linguistics --- the weight of their composite knowledge is the ``hypotenuse''
of the two components. They are not more knowledgeable than a specialist in
either field (along that field's axis), but their knowledge occupies a unique
position in meaning space that neither specialist's knowledge can reach.

The Pythagorean structure also explains why interdisciplinary work is
simultaneously valuable and difficult. It is valuable because orthogonal
combination reaches new regions of meaning space (the hypotenuse reaches
further from the origin than either leg). It is difficult because the
components are, by definition, unrelated: expertise in one domain provides no
leverage for learning the other. There are no ``synergies'' between orthogonal
ideas --- only independent contributions.

Contrast this with the case of aligned ideas ($\rs{a}{b} > 0$), where
composition produces synergy ($w(a \compose b)^2 > w(a)^2 + w(b)^2$), and the
case of opposing ideas ($\rs{a}{b} < 0$), where composition produces
cancellation ($w(a \compose b)^2 < w(a)^2 + w(b)^2$). Orthogonality is the
neutral case: no synergy, no cancellation, just independence.
\end{interpretation}


\section{Generalized Pythagorean Theorem}

The Pythagorean theorem extends to any finite collection of mutually orthogonal
ideas.

\begin{theorem}[Generalized Pythagorean Theorem]
\label{thm:gen-pythagoras}
If $a_1, a_2, \ldots, a_n \in \IS$ are pairwise orthogonal ($\rs{a_i}{a_j} = 0$
for all $i \neq j$), then:
\[
  w(a_1 \compose a_2 \compose \cdots \compose a_n)^2
  = \sum_{i=1}^n w(a_i)^2.
\]
\end{theorem}

\begin{proof}
By induction on $n$. The base case $n = 1$ is trivial. For the inductive step,
let $c = a_1 \compose \cdots \compose a_{n-1}$. We need $\rs{c}{a_n} = 0$. By
bilinearity:
\[
  \rs{c}{a_n} = \rs{a_1 \compose \cdots \compose a_{n-1}}{a_n}
  = \sum_{i=1}^{n-1} \rs{a_i}{a_n} = 0,
\]
since each $\rs{a_i}{a_n} = 0$ by hypothesis. Then by the two-variable
Pythagorean theorem:
\[
  w(c \compose a_n)^2 = w(c)^2 + w(a_n)^2
  = \left(\sum_{i=1}^{n-1} w(a_i)^2\right) + w(a_n)^2
  = \sum_{i=1}^n w(a_i)^2.
\]

In Lean~4:
\begin{lstlisting}
theorem gen_pythagoras (hs : ∀ i j, i ≠ j → rs (a i) (a j) = 0) :
    ‖φ (composeList a n)‖ ^ 2 = ∑ i in range n, ‖φ (a i)‖ ^ 2 := by
  induction n with
  | zero => simp [composeList, embed_void, norm_zero]
  | succ n ih => ...
\end{lstlisting}
\end{proof}

\begin{interpretation}[The Dimensionality of Knowledge]
\label{interp:dimensionality}
The generalized Pythagorean theorem implies that the weight of a composition
of mutually orthogonal ideas grows as $\sqrt{n}$ (if each idea has unit weight),
not as $n$. This is a \emph{curse of dimensionality} for knowledge: adding a
new orthogonal field to your expertise increases your total weight, but with
diminishing returns. The marginal contribution of each new field decreases as
$1/\sqrt{n}$.

This may help explain why genuine polymaths are rare. The weight gain from
mastering a second unrelated field (from $1$ to $\sqrt{2} \approx 1.41$) is
substantial: a 41\% increase. But the weight gain from a tenth unrelated field
(from $3$ to $\sqrt{10} \approx 3.16$) is only 5\%. Specialization offers
better returns: deepening one's existing knowledge (increasing $w(a_1)$)
increases the total weight by the full amount of the increase, with no
diminishing returns.

The exception is when the fields are \emph{not} orthogonal: when they have
positive resonance, the synergy term $2\rs{a_i}{a_j}$ boosts the total weight
beyond the Pythagorean prediction. This is why \emph{adjacent} fields combine
more productively than \emph{distant} ones: the resonance terms provide
additional weight that pure orthogonality cannot.
\end{interpretation}


\section{Orthogonal Decomposition of Ideas}

Given any idea $a$ and a non-void idea $b$, we can decompose $a$'s embedding
into components parallel and perpendicular to $b$.

\begin{theorem}[Orthogonal Decomposition]
\label{thm:orthogonal-decomposition}
For any ideas $a, b \in \IS$ with $b \not\equiv_\varphi \void$:
\[
  \embed{a} = \underbrace{\frac{\rs{a}{b}}{\rs{b}{b}} \cdot
  \embed{b}}_{\text{parallel component}} +
  \underbrace{\left(\embed{a} - \frac{\rs{a}{b}}{\rs{b}{b}} \cdot
  \embed{b}\right)}_{\text{perpendicular component}}.
\]
The parallel component has norm $\frac{|\rs{a}{b}|}{\sqrt{\rs{b}{b}}}$,
and the two components are orthogonal.
\end{theorem}

\begin{proof}
This is the standard orthogonal projection. Let $\alpha =
\frac{\rs{a}{b}}{\rs{b}{b}} = \frac{\ip{\embed{a}}{\embed{b}}}
{\nrm{\embed{b}}^2}$ and let $v_\parallel = \alpha \embed{b}$, $v_\perp =
\embed{a} - \alpha \embed{b}$. Then $\embed{a} = v_\parallel + v_\perp$ and:
\[
  \ip{v_\parallel}{v_\perp} = \alpha \ip{\embed{b}}{\embed{a} - \alpha \embed{b}}
  = \alpha \left(\ip{\embed{b}}{\embed{a}} - \alpha \nrm{\embed{b}}^2\right)
  = \alpha \left(\rs{b}{a} - \frac{\rs{a}{b}}{\rs{b}{b}} \cdot \rs{b}{b}\right)
  = 0.
\]
The norm of $v_\parallel$ is $|\alpha| \cdot \nrm{\embed{b}} =
\frac{|\rs{a}{b}|}{\nrm{\embed{b}}}$.

In Lean~4:
\begin{lstlisting}
theorem orthogonal_decomposition (hb : φ b ≠ 0) :
    φ a = (rs a b / rs b b) • φ b +
          (φ a - (rs a b / rs b b) • φ b) := by
  ring
\end{lstlisting}
\end{proof}

\begin{interpretation}[Aspect and Novelty]
\label{interp:aspect}
The orthogonal decomposition of idea $a$ with respect to idea $b$ separates $a$
into two components:

\begin{enumerate}
\item The \textbf{aspect of $a$ relative to $b$}: the part of $a$'s meaning
that lies in the direction of $b$. This is the portion of $a$ that ``looks
like'' $b$ --- that can be understood from $b$'s perspective.

\item The \textbf{novelty of $a$ relative to $b$}: the part of $a$'s meaning
that is orthogonal to $b$. This is the portion of $a$ that is genuinely new
from $b$'s perspective --- that cannot be captured by $b$ at all.
\end{enumerate}

In hermeneutic terms, this decomposition formalizes Gadamer's notion of
\emph{prejudice} (Vorurteil). When we encounter a new idea $a$, we inevitably
understand it through the lens of our existing ideas (here represented by $b$).
The ``parallel component'' is what we understand immediately: the aspect of $a$
that resonates with what we already know. The ``perpendicular component'' is
what is genuinely new: the aspect of $a$ that our existing knowledge cannot
accommodate, and that requires genuine learning to grasp.

Gadamer argued that prejudice is not merely a bias to be overcome but a
\emph{necessary condition} for understanding: we can only understand the new
by relating it to the old. The orthogonal decomposition makes this precise:
the parallel component \emph{is} the understanding afforded by prejudice,
and it is necessarily present (unless $a$ is completely orthogonal to our
existing knowledge, in which case we literally cannot understand any of it).

The ratio $\frac{|\rs{a}{b}|^2}{\rs{a}{a} \cdot \rs{b}{b}}$ --- the squared
cosine similarity --- measures the fraction of $a$'s weight that is accessible
from $b$'s standpoint. When this ratio is close to 1, almost all of $a$ is
accessible (the ``prejudice'' captures almost everything). When it is close
to 0, almost nothing is accessible (the idea is almost entirely novel).
\end{interpretation}


\section{Orthogonal Complements and Language Games}

\begin{definition}[Orthogonal Complement]
\label{def:orthogonal-complement}
For a subset $S \subseteq \IS$, the \textbf{orthogonal complement} is:
\[
  S^\perp := \{a \in \IS : \rs{a}{s} = 0 \text{ for all } s \in S\}.
\]
\end{definition}

\begin{proposition}[Properties of Orthogonal Complements]
\label{prop:orth-comp}
\begin{enumerate}
\item $\void \in S^\perp$ for any $S$.
\item If $a, b \in S^\perp$, then $a \compose b \in S^\perp$.
\item $S^\perp$ is a submonoid of $(\IS, \compose, \void)$.
\item $S \cap S^\perp \subseteq \{a : \embed{a} = 0\}$ (only void-synonymous
ideas can be both in $S$ and $S^\perp$).
\end{enumerate}
\end{proposition}

\begin{proof}
(1) $\rs{\void}{s} = 0$ for all $s$.

(2) $\rs{a \compose b}{s} = \rs{a}{s} + \rs{b}{s} = 0 + 0 = 0$ by bilinearity.

(3) Follows from (1) and (2).

(4) If $a \in S \cap S^\perp$, then $a \in S$ and $\rs{a}{s} = 0$ for all
$s \in S$. In particular, $\rs{a}{a} = 0$, so $\embed{a} = 0$.

In Lean~4:
\begin{lstlisting}
theorem orth_comp_submonoid (S : Set Idea)
    (ha : a ∈ S.orthComp) (hb : b ∈ S.orthComp) :
    a ∘ b ∈ S.orthComp := by
  intro s hs
  simp [rs_compose_left, ha s hs, hb s hs]
\end{lstlisting}
\end{proof}

\begin{interpretation}[Wittgenstein's Language Games as Orthogonal Subspaces]
\label{interp:language-games}
In the \emph{Philosophical Investigations}, Wittgenstein introduced the concept
of ``language games'' (Sprachspiele): distinct forms of linguistic activity,
each with its own rules, vocabulary, and criteria of success. The language game
of giving orders is different from the language game of telling jokes, which is
different from the language game of doing physics.

In IDT~v2, language games correspond to \emph{orthogonal subspaces} of the
Hilbert space. The ideas belonging to one language game (say, chess strategy)
form a subset $S_{\text{chess}}$ of $\IS$, and the span of their embeddings
defines a subspace of $\Hspace$. If the chess-strategy subspace is orthogonal
to the cooking-recipe subspace, then:

\begin{enumerate}
\item Ideas within each game resonate with each other but not with ideas in the
other game.
\item Composing an idea from each game produces a composite whose weight is
the Pythagorean combination (no interaction term).
\item Expertise in one game provides no leverage for the other (the parallel
component of any chess idea in the cooking direction is zero).
\end{enumerate}

This is a formal model of Wittgenstein's insight that language games are, in a
deep sense, \emph{independent}: mastery of one does not transfer to another.
But the Hilbert space model goes further than Wittgenstein. It allows for
language games that are \emph{partially} overlapping (non-zero but small
resonance), \emph{aligned} (positive resonance), or \emph{opposed} (negative
resonance). The discrete, all-or-nothing character of Wittgenstein's language
games is replaced by a continuous spectrum of inter-game relationships.

Moreover, the orthogonal complement $S^\perp$ gives us a precise notion of
``everything that has nothing to do with $S$.'' The complement of the chess
language game is the submonoid of ideas that are completely irrelevant to
chess. This is a much larger set than the complement of any particular
language game, since most ideas have nothing to do with chess.
\end{interpretation}


\section{The Hegelian Dialectic in Hilbert Space}

We have carefully distinguished orthogonality (zero resonance) from opposition
(negative resonance). Let us now develop the theory of opposition, connecting
it to Hegel's dialectic.

\begin{definition}[Dialectical Pair]
\label{def:dialectical}
Two non-void ideas $a, b \in \IS$ form a \textbf{dialectical pair} if
$\rs{a}{b} < 0$ and $\embed{a}$, $\embed{b}$ span a two-dimensional subspace
of $\Hspace$.
\end{definition}

\begin{theorem}[Weight of Dialectical Composition]
\label{thm:dialectical-weight}
If $(a, b)$ is a dialectical pair with $\hat{r}(a, b) = \cos\theta$ for
$\theta \in (\pi/2, \pi]$, then:
\[
  w(a \compose b)^2 = w(a)^2 + w(b)^2 + 2w(a)w(b)\cos\theta.
\]
In particular:
\begin{enumerate}
\item $w(a \compose b) < \sqrt{w(a)^2 + w(b)^2}$ (less than the Pythagorean
prediction);
\item If $w(a) = w(b)$ and $\theta = \pi$ (perfect opposition), then
$w(a \compose b) = 0$ (complete cancellation).
\end{enumerate}
\end{theorem}

\begin{proof}
This is the norm expansion formula (Theorem~\ref{thm:norm-expansion}) with the
substitution $\rs{a}{b} = w(a)w(b)\cos\theta$.

(1) Since $\cos\theta < 0$, we have $2w(a)w(b)\cos\theta < 0$, so
$w(a \compose b)^2 < w(a)^2 + w(b)^2$.

(2) If $w(a) = w(b)$ and $\theta = \pi$, then $\cos\pi = -1$ and:
\[
  w(a \compose b)^2 = w^2 + w^2 - 2w^2 = 0.
\]
\end{proof}

\begin{interpretation}[Thesis, Antithesis, and the Void]
\label{interp:hegel}
Hegel's dialectic proceeds in three stages: thesis, antithesis, synthesis.
The thesis proposes; the antithesis opposes; the synthesis transcends both by
incorporating the ``truth'' of each into a higher unity.

In the linear Hilbert space model, thesis and antithesis (a dialectical pair)
compose to produce something with \emph{less} weight than either alone ---
partial cancellation. In the extreme case of perfect opposition, the
composition is the void: thesis $+$ antithesis $=$ silence.

This seems to contradict Hegel: where is the synthesis? The answer is that
\emph{the synthesis is not a linear operation}. Hegelian synthesis does not
merely add thesis and antithesis; it \emph{sublates} (aufhebt) them --- lifting
both to a higher level of understanding. In Hilbert space terms, the synthesis
would be a vector that is not in the span of the thesis and antithesis, but in a
genuinely new direction. The synthesis requires moving to a dimension that
neither the thesis nor the antithesis occupied.

This means that the Hilbert space model captures the \emph{destructive} phase
of dialectic (cancellation) but not the \emph{creative} phase (synthesis). The
creative phase requires a non-linear operation --- a ``jump'' to a new dimension
--- that goes beyond the additive axiom. This is a precise mathematical diagnosis
of the limitation of the linear model, and it points toward non-linear extensions.

One possible extension: define synthesis as
$\text{synth}(a, b) = a \compose b \compose c$, where $c$ is a new idea
orthogonal to both $a$ and $b$ that provides the ``genuinely new'' content of
the synthesis. In this model, every dialectical resolution requires the
introduction of \emph{something new} from outside the dialectical pair ---
an insight from an orthogonal domain that resolves the tension.
\end{interpretation}

\begin{example}[Orthogonal Domains of Knowledge]
\label{ex:orthogonal-domains}
Consider three domains of knowledge, each represented by a subspace of
$\Hspace = \mathbb{R}^6$:
\begin{align*}
  \text{Mathematics} &: \text{span of } e_1, e_2 \\
  \text{Poetry} &: \text{span of } e_3, e_4 \\
  \text{Cooking} &: \text{span of } e_5, e_6
\end{align*}
These three domains are pairwise orthogonal: a mathematical idea has zero
resonance with a culinary idea. Within each domain, ideas can resonate strongly
(a theorem about prime numbers resonates with a theorem about algebraic
structures), but across domains, there is no interaction.

Now consider an idea that bridges two domains: ``the golden ratio in pastry''
might have embedding $(1, 0, 0, 0, 0, 1)$, with non-zero components in both
the Mathematics and Cooking subspaces. This bridging idea resonates with both
mathematical and culinary ideas, and its existence breaks the strict
orthogonality between the domains. The ``bridge'' opens a channel of resonance
between previously unrelated fields.

In the history of ideas, such bridges are rare and highly valued: the
application of mathematical optimization to recipe design, the use of poetic
metaphor in mathematical exposition, the deployment of game theory in literary
analysis. IDT~v2 models these bridges as vectors with non-zero projections
onto multiple orthogonal subspaces --- vectors that ``live in the cracks''
between established domains.
\end{example}

\begin{example}[The Dialectic of Thesis and Antithesis]
\label{ex:dialectic}
Let $\Hspace = \mathbb{R}^3$ and consider:
\begin{align*}
  \embed{\text{thesis}} &= (3, 0, 0), \\
  \embed{\text{antithesis}} &= (-3, 0, 0), \\
  \embed{\text{synthesis}} &= (0, 0, 4).
\end{align*}
Then:
\begin{itemize}
\item $\rs{\text{thesis}}{\text{antithesis}} = -9 < 0$ (strong opposition);
\item $\embed{\text{thesis} \compose \text{antithesis}} = (0, 0, 0)$ (complete
cancellation --- the dialectical destruction);
\item $\embed{\text{thesis} \compose \text{antithesis} \compose \text{synthesis}}
= (0, 0, 4)$ (the synthesis provides all the content of the resulting idea);
\item The synthesis is orthogonal to both thesis and antithesis: $\rs{\text{synthesis}}{\text{thesis}} = 0$, $\rs{\text{synthesis}}{\text{antithesis}} = 0$.
\end{itemize}
The synthesis does not ``combine'' thesis and antithesis --- it replaces the
void left by their cancellation with something entirely new.
\end{example}


\section{Bessel's Inequality and the Richness of Meaning Space}

We close this chapter with Bessel's inequality, which provides a fundamental
bound on how an idea distributes its weight across a collection of orthogonal
reference ideas.

\begin{theorem}[Bessel's Inequality for Ideas]
\label{thm:bessel}
Let $e_1, \ldots, e_n \in \IS$ be ideas with pairwise orthogonal embeddings
that are unit vectors ($\nrm{\embed{e_i}} = 1$ and $\rs{e_i}{e_j} = 0$ for
$i \neq j$). Then for any idea $a \in \IS$:
\[
  \sum_{i=1}^n \rs{a}{e_i}^2 \leq \rs{a}{a} = w(a)^2.
\]
\end{theorem}

\begin{proof}
This is Bessel's inequality applied to the vector $\embed{a}$ and the
orthonormal set $\{\embed{e_1}, \ldots, \embed{e_n}\}$.

Define $v = \embed{a} - \sum_{i=1}^n \rs{a}{e_i} \embed{e_i}$. Then $v$ is
orthogonal to each $\embed{e_i}$, and:
\begin{align*}
  0 \leq \nrm{v}^2
  &= \ip{\embed{a} - \sum_i \rs{a}{e_i}\embed{e_i}}
        {\embed{a} - \sum_i \rs{a}{e_i}\embed{e_i}} \\
  &= w(a)^2 - 2\sum_i \rs{a}{e_i}^2 + \sum_i \rs{a}{e_i}^2 \\
  &= w(a)^2 - \sum_i \rs{a}{e_i}^2.
\end{align*}
Rearranging: $\sum_i \rs{a}{e_i}^2 \leq w(a)^2$.
\end{proof}

\begin{interpretation}[The Weight Budget]
\label{interp:bessel}
Bessel's inequality states that the total ``resonance budget'' of an idea
across any orthonormal basis of reference ideas is bounded by the idea's total
weight. An idea of weight $w$ can distribute at most $w^2$ units of squared
resonance across all reference directions. If it resonates strongly with some
reference ideas, it must resonate weakly with others.

This is a formal expression of the principle that \emph{you cannot agree with
everyone equally}. An idea that resonates strongly with ``liberty'' and
``equality'' has less resonance budget left for ``tradition'' and ``authority''
(assuming these are orthogonal reference directions). The finite weight of any
idea limits its capacity for universal agreement.

In Durkheim's sociology, this connects to the tension between mechanical and
organic solidarity. In a society with mechanical solidarity (where all
individuals share the same ideas), every individual's idea vector is parallel
to the collective's. Bessel's inequality is tight: all the weight is
concentrated on one direction. In a society with organic solidarity (where
individuals specialize), the idea vectors point in many different orthogonal
directions. Bessel's inequality is loose: each individual's resonance with the
collective is only a fraction of their total weight.

The transition from mechanical to organic solidarity, in IDT~v2 terms, is a
transition from a low-dimensional culture (all ideas aligned) to a
high-dimensional culture (ideas spanning many orthogonal directions). The
total ``cultural weight'' increases (many specialists contribute), but the
individual resonance with the collective decreases (each specialist is only
partially aligned with the whole).
\end{interpretation}

\bigskip

This concludes Part~I of the book. We have established the Hilbert space
foundations of IDT~v2: the embedding axioms
(Chapter~\ref{ch:geometry}), resonance as inner product
(Chapter~\ref{ch:resonance}), the metric of meaning (Chapter~\ref{ch:metric}),
the independence of depth and norm (Chapter~\ref{ch:depth-norm}), and the
structure of orthogonal ideas (this chapter). In Part~II, we turn to the
\emph{game theory of meaning}: what happens when ideas are not merely objects
to be contemplated, but strategic instruments deployed by agents who seek to
influence each other?
